\documentclass[lualatex,paper=a4paper]{jlreq}
\usepackage{enumitem}
\usepackage{amsmath}
\usepackage{amssymb}
\usepackage{amsthm}
\usepackage{amsfonts}  %以上必須パッケージ
\usepackage{mathtools} %参照していない数式に番号をつけない
\usepackage{graphicx}  %図を入れる
\usepackage{multirow}  %表の列結合
\usepackage{hyperref}  %文章内にリンクを貼る
\usepackage{diffcoeff} %diff、diffpで微分記号を簡潔に書く
\usepackage{comment}   %comment環境
\usepackage{pdfpages}  %PDFページを直接挿入
\usepackage{float}     %[H]指定子でその位置に強制的に図表を配置する
\usepackage[version=4]{mhchem} %\ce{$n$H3O+}のようにして簡潔に化学式を書く
\usepackage[separate-uncertainty]{siunitx}  %SI単位を\unitと\qtyを使ってより簡便に書く
\usepackage[math-style=ISO,warnings-off={mathtools-colon,mathtools-overbracket}]{unicode-math} %Unicode文字を数式に書けるようにする
\usepackage{newunicodechar} %Unicode文字を命令化できるようにする
\usepackage{zref-clever}    %\crefで参照の種類に応じて「式」「図」「表」などと自動的に出力してくれる
\usepackage{listings}       %ソースコードを簡単に書けるようにする
%%%%%%%%%%%%%%%%%%%%%%%%%%%%%%%%%%%%%%%%%%%%%%%%%
\mathtoolsset{showonlyrefs,showmanualtags} %参照されていない数式に番号をつけない
%%\newunicodechar{、}{，}  %、を，に自動置換
%%\newunicodechar{。}{．}  %。を．に自動置換
\NewDocumentCommand\、{}{\string、} %\、で読点をそのまま出力
\NewDocumentCommand\。{}{\string。} %\。で句点をそのまま出力
\NewDocumentCommand\degC{}{\ensuremath{^\circ\symup{C}}} %摂氏温度記号
\NewDocumentCommand\abs{m}{\left|#1\right|} %絶対値の記号
\renewcommand{\lstlistingname}{コード} %listingsの環境名を「コード」に変更
%%%%%%%%%%%%%%%%%%%%%%%%%%%%%%%%%%%%%%%%%%%%%%%%%
\setmathfont{NewComputerModernMath}[
    CharacterVariant = {1} %∅(空集合)の記号を一般的なものに変更
]
\setmathfont{NewComputerModernMath}[
    range = {scr, bfscr},  %花文字のフォントを変更する、これがないと普通の筆記体として表示される
    StylisticSet = 1
]
% zref-cleverの日本語設定
\zcDeclareLanguage[allcaps]{japanese}
\zcLanguageSetup{japanese}{
  namesep   = {\nobreak} ,
  pairsep   = {と} ,        % 複数個の参照をつなぐときの接続詞
  listsep   = {、} ,        % 複数個の参照をつなぐときの接続詞（3個以上）
  lastsep   = {、および} ,  % 複数個の参照をつなぐときの最後の接続詞（3個以上）
  tpairsep  = {および} ,    % 複数個の異なる種類の参照をつなぐときの接続詞（2個）
  tlistsep  = {、} ,        % 複数個の異なる種類の参照をつなぐときの接続詞（3個以上）
  tlastsep  = {、および} ,  % 複数個の異なる種類の参照をつなぐときの最後の接続詞（3個以上）
  notesep   = {：} ,        % 参照の後に注釈をつけるときの区切り
  rangesep  = {から} ,      % 範囲指定のときの区切り
  type = section ,          % 節の引用に関する設定
    Name-sg = 第 ,
    Name-pl = 第 ,
    refbounds = {,,節,} ,
  type = subsection ,       % 小節の引用に関する設定
    Name-sg = 第 ,
    Name-pl = 第 ,
    refbounds = {,,小節,} ,
  type = subsubsection ,    % 小小節の引用に関する設定
    Name-sg = 第 ,
    Name-pl = 第 ,
    refbounds = {,,小小節,} ,
  type = paragraph ,
    Name-sg = 第 ,
    Name-pl = 第 ,
    refbounds = {,,段落,} ,
  type = appendix ,
    Name-sg = 付録 ,
    Name-pl = 付録 ,
  type = page ,             % ページの場合
    Name-sg = p. ,
    Name-pl = pp. ,
    namesep = 〜 ,
    rangesep = \textendash ,
    rangetopair = false ,
  type = figure ,
    Name-sg = 図 ,
    Name-pl = 図 ,
  type = table ,
    Name-sg = 表 ,
    Name-pl = 表 ,
  type = item ,
    Name-sg = 項目 ,
    Name-pl = 項目 ,
  type = footnote ,
    Name-sg = 脚注 ,
    Name-pl = 脚注 ,
  type = endnote ,
    Name-sg = 後注 ,
    Name-pl = 後注 ,
  type = note ,
    Name-sg = 注記 ,
    Name-pl = 注記 ,
  type = equation ,
    Name-sg = 式 ,
    Name-pl = 式 ,
    refbounds = {,(,),} ,
  type = listing ,
    Name-sg = コード ,
    Name-pl = コード ,
}
\zcsetup{ lang = japanese, nameinlink = true } % 日本語設定を有効化
%%%%%%%%%%%%%%%%%%%%%%%%%%%%%%%%%%%%%%%%%%%%%%%%%
\lstset{%
  basicstyle={\ttfamily}, % コード全体の書体
  identifierstyle={\normalsize}, % キーワード以外のコードの書体
  keywordstyle={\small\bfseries\color[rgb]{0.000, 0.000, 1.000}}, % キーワード(ifなど)の書体
  backgroundcolor={\color[gray]{0.98}}, % 背景色
  stringstyle={\small\ttfamily\color[rgb]{0.639, 0.082, 0.082}}, %文字列の書体
  commentstyle={\small\ttfamily\color[rgb]{0.000, 0.502, 0.000}}, %コメントの書体
  frame={tb}, % コードの上下に線を引く。lrを足すと左右にも線が引かれる
  breaklines=true, % 行が長くなったときに自動で改行する
  showstringspaces=false, % 文字列中のスペースを可視化(␣で表記)しない
  columns=fixed, % 等角で文字を配置する
  basewidth=0.5em, % 等幅フォントの文字幅を0.5emに設定
  numbers=left, % 行番号を左に表示
  xrightmargin=0\zw, %右側の余白を全角幅0に設定
  xleftmargin=3\zw, %左側の余白(行番号表示部分)を全角幅3に設定
  numberstyle={\scriptsize}, % 行番号の書体
  stepnumber=1, % 行番号を1行ごとに表示
  numbersep=1\zw, % 行番号とコードの間のスペースを全角幅1に設定
  lineskip=-0.5ex % 行間を詰める
}
%%%%%%%%%%%%%%%%%%%%%%%%%%%%%%%%%%%%%%%%%%%%%%%%%
% jlreqの設定を変更するためのコマンド
% 詳細: https://tug.org/texlive//Contents/live/texmf-dist/doc/latex/jlreq/jlreq-ja.html
\jlreqsetup{
    appendix_counter={ %付録に入ったときの処理
        section={
            value=0,
            the={\Alph{section}} %付番を戻してアルファベット(A, B, C...)に変更
        },
        table={
            value=0,
            the={\Alph{section}\arabic{table}} %表番号を付録のものに変更
        },
        figure={
            value=0,
            the={\Alph{section}\arabic{figure}} %図番号を付録のものに変更
        }
    },
    appendix_heading={
        section={
            font=\normalsize\bfseries,
            label_format={付録\thesection:} %付録の表記を変更
        },
        subsection={
            font=\normalsize\bfseries
        },
        subsubsection={
            font=\normalsize\bfseries
        }
    },
    caption_font=\mcfamily, %図表の題を明朝体で出力する
    caption_label_font=\mcfamily, %図表番号を明朝体にする
    caption_label_format={#1:}, %「図1:」のようにする
    caption_after_label_space=0.5\zw %:の後に全角幅半分のスペースを空ける
}
\ModifyHeading{section}{ %見出しの表示を変更
    font=\normalsize\bfseries, %見出しのサイズを変更
    lines=2, %見出しの行取りを変更
    after_lines=0
}
\ModifyHeading{subsection}{ %小見出しの表示を変更
    font=\normalsize\bfseries,
    lines=2,
    after_lines=0
}
\ModifyHeading{subsubsection}{ %小小見出しの表示を変更
    font=\normalsize\bfseries,
    lines=2,
    after_lines=0
}

\title{ゼロからpythonで行列の簡約化} % タイトルを入れる
\author{MMA平部員 B1 Tenpura} % 学籍番号と名前を入れる
\date{\today}

\begin{comment}
このテンプレートはhydrogen(https://hydrogen1.dev)によって作成されました。改変や再頒布、どのように扱っても問題ありません。
元々のテンプレートはこちら(https://gitlab.mma.club.uec.ac.jp/hydrogen/latex-template)にあります。
邪魔であれば、このコメントの部分(\begin{comment}からend{comment}まで)を消してください。

また、このテンプレートはLuaLaTeXを対象に作られています。
改変していないテンプレートから正常にPDFを生成できない場合、以下のURLからTeXLiveをインストールしてください。
https://mirror.ctan.org/systems/texlive/tlnet/install-tl-windows.exe
\end{comment}

\begin{document}
%%% --- ここから書き始める
%表紙が不要な場合は下の行頭の%を取り除く
%\begin{comment}
%\begin{titlepage}
%    \includepdf{cover_sample.pdf} %表紙のPDFファイルを挿入する。
%\end{titlepage}
%%\end{comment}
%表紙が不要な場合は上の行頭の%を取り除く

\maketitle %タイトルの出力
\section{はじめに}
こんにちは。B1平部員のTenpuraです。この記事では夏休みに重い腰を上げて知識ゼロからPythonで行列式を書いてみた話をします。
動機はAnkiのアドオン作りでPythonわからんことによる萎え落ちを減らすことです。あと線形代数で追試くらったことです。雪辱を果たしに。

\section{行列とは}
行列は縦横に数を並べたものです。例えば
$\begin{cases}
        x + 2y - z = 2  \\
        2x - y + 3z = 9 \\
        3x + y + z = 8
    \end{cases}$
を
$\begin{bmatrix}
        1 & 2  & -1 & 2 \\
        2 & -1 & 3  & 9 \\
        3 & 1  & 1  & 8
    \end{bmatrix}$

のように書けます。シンプルですね。

\section{簡約化とは}
簡約化は、連立方程式を解く際に行列の値を都合よくする操作です。実際中学からの連立方程式の解き方まんまで、式同士を足し引きします。
例えば、３つの式を上から①②③として、

$$\begin{bmatrix}
        1 & 2  & -1 & 2 \\
        2 & -1 & 3  & 9 \\
        3 & 1  & 1  & 8
    \end{bmatrix}
    \xrightarrow[\text{③}-3\times\text{①}]{\text{②}-2\times\text{①}}
    \begin{bmatrix}
        1 & 2  & -1 & 2 \\
        0 & -5 & 5  & 5 \\
        0 & -5 & 4  & 2
    \end{bmatrix}
    \xrightarrow{\text{②}\times(-\frac{1}{5})}
    \begin{bmatrix}
        1 & 2  & -1 & 2  \\
        0 & 1  & -1 & -1 \\
        0 & -5 & 4  & 2
    \end{bmatrix}$$
$$\xrightarrow[\text{③}+5\times\text{②}]{\text{①}-2\times\text{②}}
    \begin{bmatrix}
        1 & 0 & 1  & 4  \\
        0 & 1 & -1 & -1 \\
        0 & 0 & -1 & -3
    \end{bmatrix}
    \xrightarrow{\text{③}\times(-1)}
    \begin{bmatrix}
        1 & 0 & 1  & 4  \\
        0 & 1 & -1 & -1 \\
        0 & 0 & 1  & 3
    \end{bmatrix}
    \xrightarrow[\text{②}+\text{③}]{\text{①}-\text{③}}
    \begin{bmatrix}
        1 & 0 & 0 & 1 \\
        0 & 1 & 0 & 2 \\
        0 & 0 & 1 & 3
    \end{bmatrix}$$
と変形します（行基本変形と言います）。最後の行列は

$$\begin{cases}
        x + 0 + 0 = 1 \\
        0 + y + 0 = 2 \\
        0 + 0 + z = 3
    \end{cases}$$

となり、良い感じに文字の値が分かります。
この操作を言語化すると、\\
\hrulefill
\begin{verbatim}
1.すべての成分が0でない列のうち、一番左の列に注目する。
2.注目した列の一番上の数が0のときは、0でない行と入れ替える。
3.第1行を何倍かして主成分をつくる。
4.第1行を何倍かして他の行に加えて、主成分の上下にある成分をすべて0にする。
5.第1行を忘れ、残った行で1～4を繰り返す。
\end{verbatim}
\hrulefill\\
となります。\footnote{培風館 線形代数学講義[増補版] 木田雅成 著 P40 Ⅲ.7より引用}
我々B1は4月の不安な時期にこれを頭に叩き込んだと思います。\\
個人的な話ですが4行4列よりサイズが大きくなると必ず計算ミスをします。どうにかしたい。

\section{計画}\label{plan}
今回はこれをPythonで実装します。しかし、上記の説明のまま実装しようとすると割と苦労します。例えば5.で繰り返しの操作がありますが、コードで書くには行数を文字で置いたりして一般化する必要があります。そこで、以下のように考えてみます。\\
\hrulefill
\begin{verbatim}
行列を縦I、横Jの大きさとする。はじめ、k=1とする。
1.左から右の列へ、各列でk行目から順に0でない数を探す。はじめに見つかる数をnとする。
2.nがある行をk行目と入れ替える。
3.k行目をnで割る。
4.k行目以外の行からk行目を何倍かして引き、nと同じ列のn以外の数を0にする。
5.kに1を加え、1.からやり直す。1.でk>Iとなるか、J列目を超えたとき、計算を終了する。
\end{verbatim}
\hrulefill\\
使いやすい文字をいくつか入れてみました。いっぺんに実装を考えるときついですが、順番ごとに考えるとなんかいけそうな気がします。
\newpage
\section{Pythonのお約束}
知識ゼロからPythonを書くにあたって、Pythonのお約束を知っておこう。Pythonを使う準備や基本構文はPython公式入門講座\footnote{https://www.python.jp/train/index.html}が簡単に教えてくれます。コンリテの数倍分かりやすいので気楽にやりましょう。
この章では補足的に、行列の書き方と「おまじない」を紹介します。

\subsection{リスト内包表記}
Pythonで行列を扱う際、リスト内包表記が便利です。リスト\footnote{複数の要素をまとめて扱う形式。日常の「リスト」と同じ。}の中にリストを入れる書き方です。
例えば$\begin{bmatrix}
        1 & 2  & -1 & 2 \\
        2 & -1 & 3  & 9 \\
        3 & 1  & 1  & 8
    \end{bmatrix}$
を$\begin{aligned}
        A =[[1,2,-1,2], \\
        [2,-1,3,9],     \\
        [3, 1,1,8]]     \\
    \end{aligned}$
と書きます。\\改行は無くてもよく、一行に直すと
A =[[1,2,-1,2],[2,-1,3,9],[3,1,1,8]]となります。\\
これを使い倒すなら行列の大きさや一つ一つの数が変数に欲しい所です。以下に実装に使うものをまとめます。

\begin{table}[htbp]
    \centering
    \caption{リスト内包表記の行列}
    \begin{tabular}{cc}
        \hline
        $a_{ij}$(i行j列目の数) & a[i-1][j-1] \\
        行数                & len(a)      \\
        列数                & len(a[0])   \\ \hline
    \end{tabular}
\end{table}
注意！！Pythonのリストは順番を0から数えます。例えば$a_{23}$はa[1][2]です。この数え方は大体のプログラミング言語で共通し、度々間違えてバグの温床となります。私は現在進行形で間違えます。記事のどっかには絶対間違ってますが許して～～～

\subsection{numpy,simpy}
他にはnumpyやsimpyというライブラリ\footnote{便利機能が集まっている道具箱}で行列を扱えます。こちらは簡約化、行列式の計算、単位行列、ゼロ行列などのプリセットが入っていて、行列でやりたい操作は大体何でもできます。例えば今回実装する簡約化はAを行列としてA.rref()で終わります。これで終わるのはさすがにおもんなさすぎるので、今回はリスト内包表記で頑張ります。
\newpage
\subsection{おまじない}
我々初心者がわけもわからずファイル先頭に入れさせられる謎の文字列こと、おまじない。正体は\LaTeX の\verb|\usepackage{}| と同じく、使うライブラリを書いておく部分です。先ほどのnumpyなどマジで便利なものがあるのでよく弄りましょう！！\\
ただ、今回は実装を楽しみたいので以下の基本的なものだけを入れます。
\begin{lstlisting}[language=Python,caption={便利なおまじない},label=code-sample]
import re
from decimal import Decimal
from fractions import Fraction
\end{lstlisting}
「re」は正規表現、「Decimal」は正確な小数計算、「Fraction」は分数を扱います。どれも便利。

\section{実装}
ようやく実装に入ります。\ref{plan}章に従ってコードを書いていきます。\\
テストの際おまじないを忘れないでください。また、使う行列はA =[[1,2,-1,2],[2,-1,3,9],[3,1,1,8]]として、以下のように実装していきます。

\begin{lstlisting}[language=Python,caption={概形},label=code-sample]
import re
from decimal import Decimal
from fractions import Fraction

A =[[1,2,-1,2],[2,-1,3,9],[3,1,1,8]]

def simplify(a):
    k = 0
    while k < len(a):
        [1つ目の処理]
        [2つ目の処理]
        [3つ目の処理]
        [4つ目の処理]
    return a 

\end{lstlisting}


\begin{lstlisting}[language=Python,caption={左から右の列へ、各列でk行目から順に0でない数を探す。はじめに見つかる数をnとする。},label=code-sample]
def simplify(a):
    k = 0
    while k < len(a):
        i = process_row
        j = 0
        while j < len(a[0]) and a[i][j] == 0:
            if i == len(a)-1:
                i = k
                j = j + 1
            else:
                i = i + 1
        if j >= len(a[0]):
            break
        print(f"a{i+1,j+1}に注目")
        tmp_i = i
        tmp_j = j
\end{lstlisting}


プログラムのソースコードをレポートに載せたいときは、listingsパッケージを利用して\verb|lstlisting|環境で囲むと良いだろう。例えば、\zcref{code-sample}のように書くとPythonのコードがきれいに表示される。

\begin{lstlisting}[language=Python,caption={便利なおまじない},label=code-sample]
def hello():
    print("Hello, World!")

if __name__ == "__main__":
    hello() # この行はコメントです
\end{lstlisting}

\subsection{参考文献の明示}

レポートを書く上では他の文献から情報を得ることも多いであろう。こういう時は、thebibliography環境に参考文献として書籍やネットの記事を並べて出典を明確にするべきである。
ネット上の情報の引用の際には\verb|\url|マクロを利用して\url{https://www.example.com/}のように書くと良いだろう\cite{texbook}。

参考文献を書くときは、\verb|\bibitem|マクロを用いて文献を列挙する。このマクロに指定したラベルを\verb|\cite|マクロで参照することで、文献番号を出力できる。

\begin{thebibliography}{99}
    \bibitem{texbook} 奥村晴彦，黒木裕介， [改訂第8版] LaTeX2ε美文書作成入門，技術評論社，2020.
\end{thebibliography}

\appendix

\section{付録の書き方}

本文に書くべきではない細かい計算や、導出は付録として別個にまとめることが多い。付録に入る際は\verb|\appendix|マクロを用いる。その後、附番はアルファベットになり表番号や図番号の表記も変更される。
なお、\verb|\appendix|は1度書いた後はもう書く必要はない。

%%% --- ここまで
\end{document}